\chapter{Reproducibilidad}
\label{cap:apendice}

\section{Estructura del proyecto}

\begin{verbatim}
PROYECTO1_PLANIFICA/
|-- data/raw/                 dataset oficial (CSV) + leaderboard publico
|-- codigos/
|   |-- ejecutados/           4 notebooks reproducidos (.ipynb + .py)
|   |-- descartados/          8 notebooks analizados sin ejecutar
|   `-- _RESUMEN_TECNICAS.md  inventario de tecnicas por notebook
|-- src/                      benchmarks reproducibles
|   |-- bench_01_inversion.py baseline XGB CPU
|   |-- bench_02_xhlulu.py    baseline XGB GPU
|   |-- bench_03_nroman.py    LightGBM + FE (TimeSeriesSplit)
|   |-- bench_04_cdeotte.py   XGB d12 con/sin magic UID
|   |-- download_notebooks.py descarga via API publica de Kaggle
|   |-- extract_notebooks.py  exportacion .ipynb -> .py
|   `-- make_ranking.py       consolidacion de resultados
|-- results/                  CSV y logs de cada benchmark + ranking
`-- informe_latex/            este informe (tectonic main.tex)
\end{verbatim}

\section{Entorno}

\begin{itemize}
  \item Linux; Python 3.12 en entorno \texttt{uv}; pandas 3.0, XGBoost 3.4,
        LightGBM 4.x, scikit-learn.
  \item CPU 16 núcleos; NVIDIA RTX 2060 6\,GB (driver 610).
  \item Datos: \texttt{kaggle competitions download -c ieee-fraud-detection}
        (requiere aceptar las reglas de la competencia).
\end{itemize}

\section{Ejecución de los experimentos}

\begin{verbatim}
uv venv --python 3.12 .venv
uv pip install pandas numpy scikit-learn lightgbm xgboost
uv pip install matplotlib
.venv/bin/python src/bench_01_inversion.py     # AUC 0.9191, 112 s
.venv/bin/python src/bench_02_xhlulu.py        # AUC 0.9227,  25 s (GPU)
.venv/bin/python src/bench_03_nroman.py        # AUC 0.9200, 462 s
.venv/bin/python src/bench_04_cdeotte.py       # 0.9375 -> 0.9479 (GPU)
.venv/bin/python informe_latex/scripts/generar_graficos.py
tectonic informe_latex/main.tex
\end{verbatim}

\section{Adaptaciones de versiones (2019 $\to$ 2026)}

\begin{itemize}
  \item \textbf{XGBoost $\ge$2}: \texttt{tree\_method='gpu\_hist'} se reemplaza por
        \texttt{tree\_method='hist', device='cuda'}; \texttt{early\_stopping\_rounds} y
        \texttt{eval\_metric} pasan al constructor del estimador.
  \item \textbf{LightGBM $\ge$4}: los argumentos \texttt{early\_stopping\_rounds} y
        \texttt{verbose\_eval} de \texttt{lgb.train} se sustituyen por los
        \emph{callbacks} \texttt{lgb.early\_stopping()} y \texttt{lgb.log\_evaluation()}.
  \item \textbf{pandas 3}: las cadenas leídas de CSV adoptan el dtype \texttt{str} (no
        \texttt{object}); las condiciones \texttt{dtype=='object'} deben reemplazarse por
        \texttt{pd.api.types.is\_numeric\_dtype(...)} negado. \texttt{LabelEncoder} se
        sustituye de forma segura por \texttt{pd.factorize(sort=True)} cuando hay tipos mixtos.
  \item \textbf{RAPIDS cuDF} (notebook GPU de Deotte): requiere $>$6\,GB de VRAM; no ejecutable
        en la RTX 2060, se analizó cualitativamente.
\end{itemize}

\section{Protocolo ético}

El estudio se limita a material público con licencia abierta (Apache 2.0 en los notebooks
citados). No se realizó ninguna submisión a la competencia ni uso de servicios externos con los
datos; toda la evidencia numérica proviene de ejecuciones locales.
